\documentclass[11pt]{article}

\usepackage[margin=1in]{geometry}
\usepackage{amsmath,amssymb,amsthm}
\usepackage{bm}
\usepackage{hyperref}
\usepackage{float}
\usepackage{graphicx}
\usepackage{booktabs}
\usepackage{subcaption}
\usepackage{verbatim}
\usepackage{longtable}
\usepackage{fontspec}

\setmainfont{Times New Roman}
\setmonofont{Menlo}

\graphicspath{{./}{docs/models/}{docs/models/images/}}

\setlength{\parindent}{0pt}
\setlength{\parskip}{0.7\baselineskip}

\title{Report: Single-Model PyMC Collapsed Ridge RT Baseline (Seen Groups)}
\author{Bioinformatics Team}
\date{\today}

\begin{document}

\maketitle

\section{Executive summary}

This document summarises our current RT modelling findings for the \textbf{seen-group} setting
(predicting new runs for groups $g = (\texttt{species\_cluster}, \texttt{comp\_id})$ present in training),
and compares:
\begin{itemize}
  \item the existing \textbf{sklearn per-group ridge} artifact (Stage-1 coefficient summaries), and
  \item a \textbf{PyMC collapsed ridge} formulation trained from per-group sufficient statistics.
\end{itemize}

We initially explored a PyMC ``single model'' direction to enable extensions (hierarchical pooling and coherent uncertainty), but the
results below show that for cap100 data the simplest ridge baselines already dominate for point prediction.

\paragraph{Key outcome (cap100 $\rightarrow$ realtest).}
Training on full cap100 and evaluating on realtest (seen groups only), we can improve beyond the linear ridge baseline by adding a
lightweight nonlinearity in the \texttt{IS*}/\texttt{RS*} anchor predictors:
\textbf{degree-2 (squares + interactions) on anchors only}, with ridge shrinkage.
This yields the best RMSE we have observed so far:
lib208 $\approx 0.00891$ and lib209 $\approx 0.00800$.

\paragraph{Why RMSE matches ridge (when $\lambda$ is fixed).}
With the Gaussian prior $\bm{w}_g \sim \mathcal{N}(\bm{0}, \sigma_y^2 \bm{\Lambda}^{-1})$ (diagonal $\bm{\Lambda}$; optionally
two-$\lambda$ for base vs poly2 terms) and a flat intercept prior, the posterior mean of $\bm{w}_g$ conditional on
$\bm{\Lambda}$ is exactly the ridge solution:
$\hat{\bm{w}}_g = (\bm{X}_c^\top \bm{X}_c + \bm{\Lambda})^{-1}\bm{X}_c^\top \bm{y}_c$.
Therefore, when running \texttt{--lambda-mode fixed}, the PyMC model is intentionally ridge-equivalent for point predictions,
so RMSE matches sklearn ridge up to numerical precision.

\paragraph{Truth: do we need PyMC here?}
For \textbf{seen-group} RT at cap100, PyMC has not improved RMSE relative to the best sklearn coefficient-summary baselines so far.
In particular, the strongest RMSE lift we found came from a \emph{feature} change (a nonlinear anchor expansion) that does not require
PyMC. Additional PyMC pooling structures we tested (\texttt{chem\_hier} embeddings, \texttt{comp\_hier} intercept pooling, and a
per-\texttt{comp\_id} slope head) also did not improve RMSE.
The strongest practical reason to use PyMC would be if we need model extensions (e.g., different likelihoods or hierarchical pooling that
improves tail behaviour) or if we want to propagate more sources of uncertainty than the closed-form ridge approximation.

\paragraph{Uncertainty for downstream peak assignment.}
If the main requirement is per-row predictive uncertainty for peak assignment, PyMC is not strictly required: the sklearn Stage-1 artifact
already stores full coefficient covariance (\texttt{beta\_cov}) and a Bayesian noise estimate (\texttt{sigma2\_mean}), enabling an analytic
Normal predictive distribution per group. However, for the strongest RMSE setting we found (poly2 anchors), the raw per-group Normal
intervals under-cover and need a simple calibration step (e.g., a global \texttt{pred\_std\_scale} fit on a validation set).
In contrast, the PyMC collapsed ridge poly2 variant (with fixed $\lambda$) tends to produce closer-to-nominal coverage out-of-the-box via
a global $\sigma_y$ estimate. Separately, the chem\_hier PyMC variants tended to be much more conservative (coverage far above 0.95),
which may or may not be desirable downstream depending on the assignment objective.

\paragraph{Where results live.}
\begin{itemize}
  \item Figures: \texttt{docs/models/images/rt\_pymc\_collapsed\_ridge\_report/}
  \item Metric JSON/CSV: \texttt{docs/models/results/rt\_pymc\_collapsed\_ridge\_report/}
\end{itemize}

\section{Beyond seen-group RMSE: where a custom PyMC model can help}

The benchmark above answers a narrow question: \emph{for a fixed feature set, predicting RT for seen groups
$g=(\texttt{species\_cluster},\texttt{comp\_id})$ in held-out realtest runs, does additional Bayesian structure improve global RMSE?}
In this setting, ridge-equivalent linear models already saturate RMSE, and the pooling structures we tested did not help.

This does \emph{not} imply ``PyMC is useless'' for RT. It implies that to beat ridge on RMSE we must either (a) introduce additional
signal (features/conditioning information), or (b) change the functional form of the predictor. PyMC is most valuable when we want a
probabilistic model whose \emph{uncertainty} directly supports downstream peak assignment, or when we want to incorporate latent
structure that is awkward to maintain in a purely point-estimate pipeline.

\subsection{Mean function vs noise model}

It is useful to separate:
\[
  \mu(x,g) = \mathbb{E}[\texttt{rt} \mid x,g] \quad \text{(mean function)}, \qquad
  \texttt{rt} \mid x,g \sim \text{Likelihood}(\mu(x,g), \ldots) \quad \text{(noise model)}.
\]
Ridge and the collapsed Bayesian ridge baseline both assume a linear mean function in the run covariates, with Gaussian noise.
Changing the \emph{inference engine} (sklearn $\rightarrow$ PyMC) without changing $\mu(\cdot)$ typically does not move RMSE; it mainly
changes how uncertainty is represented and propagated.

\subsection{Run-level latent calibration / warping (peak-assignment use)}

In the Sally peak-assignment pipeline, predicted RT is primarily used as a \emph{gating window} (filter peaks by
$|\texttt{apex\_rt}-\texttt{expectedRT}| \le \texttt{expectedRTWindow}$; see
\texttt{external\_repos/sally/sally/analytical/curationModel.py}).
This creates an opportunity that is \emph{not} measured by our current ``new runs only'' realtest benchmark:
if we are allowed to condition on a small amount of within-run information (calibrants, confident early assignments, or a subset of
standards), we can infer a run-specific correction and tighten/shift expected RTs \emph{for that run}.

Examples of mean-function extensions:
\begin{itemize}
  \item \textbf{Per-run offset/scale:} $\mu(x,g,r) = b_g + x_r^\top w_g + a_r + s_r \,\tilde{t}_g$.
        Here $(a_r,s_r)$ are latent run parameters and $\tilde{t}_g$ is a group baseline position (e.g., learned compound intercept).
  \item \textbf{Monotone run warp:} $\mu(x,g,r) = b_g + x_r^\top w_g + f_r(\tilde{t}_g)$ with $f_r(\cdot)$ a monotone spline.
\end{itemize}
These models can be updated in (semi) real time because each run has only a small number of latent parameters; inference can be MAP +
Laplace, variational, or an analytically integrated conditional where possible. The payoff is improved downstream decisions (fewer false
peak candidates, fewer missed peaks) even if \emph{global} RT RMSE across all test rows is already near-optimal.

\subsection{Hierarchical pooling targeted at sparse regimes}

Our cap100 data has strong coverage for many groups, so additional pooling does not move \emph{global} RMSE. However, pooling can still
be valuable when focusing on sparse slices (long-tail compounds and under-filled species$\times$compound pairs), where per-group fits are
high variance. A custom PyMC model can implement adaptive partial pooling such as:
\[
  w_g = w_0 + w_{\texttt{cluster}(g)} + w_{\texttt{comp}(g)} + \delta_g,
\]
with learned shrinkage scales so well-covered groups remain effectively unpooled while sparse groups borrow strength.
To evaluate this honestly, report RMSE stratified by group support (e.g., by $n_g$ or by per-compound support), not only the global
average.

\subsection{Robust and heteroscedastic likelihoods (decision-calibrated uncertainty)}

Residuals on realtest include occasional heavy tails (mis-integrations/outliers). A Student-$t$ likelihood, mixture noise, or
heteroscedastic variance (e.g., per cluster or per run) can make fitting and predictive intervals more robust.
This may not always lower RMSE, but it can materially improve the calibration/utility of the RT window used downstream.

\subsection{Joint modelling with assignment feedback (advanced)}

The most ``Bayesian'' extension is to couple RT prediction with peak assignment: treat candidate peaks as a mixture of
``correct peak'' vs ``distractors'', use the RT model as a likelihood term, and update the run calibration and assignment confidence
iteratively as evidence accumulates. This is heavier engineering, but it targets the true downstream objective (assignment accuracy)
rather than RT RMSE alone.

\section{Code inventory (what to use)}

\paragraph{Baselines.}
\begin{itemize}
  \item Lasso (eslasso baseline): \texttt{scripts/pipelines/eval\_rt\_lasso\_baseline\_by\_species\_cluster.py}
  \item Ridge (sklearn per-group summaries): \texttt{scripts/pipelines/train\_rt\_stage1\_coeff\_summaries.py}
\end{itemize}

\paragraph{Single-model PyMC baseline.}
\begin{itemize}
  \item Train (collapsed ridge): \texttt{scripts/pipelines/train\_rt\_pymc\_collapsed\_ridge.py}
\end{itemize}

\paragraph{Shared evaluation and comparison.}
\begin{itemize}
  \item Evaluate any \texttt{Stage1CoeffSummaries}: \texttt{scripts/pipelines/eval\_rt\_coeff\_summaries\_by\_species\_cluster.py}
  \item Compare-by-cluster plot: \texttt{scripts/pipelines/compare\_rt\_models\_by\_cluster.py}
\end{itemize}

\section{Model specification}

\subsection{Data and grouping}

We train and evaluate on production RT CSVs containing:
\begin{itemize}
  \item response: $y$ = \texttt{rt},
  \item group key: $g = (\texttt{species\_cluster}, \texttt{comp\_id})$,
  \item run covariates: $\bm{x} \in \mathbb{R}^p$ from columns \texttt{IS*}, \texttt{RS*}, and optionally \texttt{ES\_*}.
\end{itemize}

Optional global feature transform (pure reparameterization for stability):
\[
  \bm{x}' = (\bm{x} - \bm{\mu}) \, \bm{R},
\]
where $\bm{\mu}$ is the global mean over training rows (optional) and $\bm{R}$ is an optional orthonormal rotation
(e.g. PCA rotation). The artifact stores $\bm{\mu}$ and $\bm{R}$ so evaluation applies the same transform.

\subsection{Generative model (per group)}

For each group $g$ with $n_g$ observations:
\[
  y_{gi} = b_g + \bm{x}_{gi}^\top \bm{w}_g + \varepsilon_{gi}, \quad
  \varepsilon_{gi} \sim \mathcal{N}(0, \sigma_y^2).
\]

We use a ridge prior on slopes:
\[
  \bm{w}_g \mid \sigma_y^2, \lambda \sim \mathcal{N}\!\left(\bm{0}, \frac{\sigma_y^2}{\lambda}\bm{I}_p\right),
\]
and treat the intercept $b_g$ as unpenalized (equivalently: an improper flat prior on $b_g$).

Hyperpriors:
\[
  \sigma_y \sim \mathrm{HalfNormal}(s_{\sigma}), \qquad
  \lambda \sim
  \begin{cases}
    \text{fixed} & (\texttt{--lambda-mode fixed}) \\
    \mathrm{LogNormal}(\mu_{\lambda}, s_{\lambda}) & (\texttt{--lambda-mode learn})
  \end{cases}
\]

\subsection{Collapsed marginal likelihood (sufficient statistics)}

Define per-group means and centered quantities:
\[
  \bar{y}_g = \frac{1}{n_g}\sum_i y_{gi}, \qquad
  \bar{\bm{x}}_g = \frac{1}{n_g}\sum_i \bm{x}_{gi},
\]
\[
  \bm{y}_{c} = \bm{y}_g - \bar{y}_g \bm{1}, \qquad
  \bm{X}_{c} = \bm{X}_g - \bm{1}\bar{\bm{x}}_g^\top.
\]

Sufficient statistics:
\[
  \bm{S}_{xx} = \bm{X}_{c}^\top \bm{X}_{c} \in \mathbb{R}^{p\times p}, \quad
  \bm{s}_{xy} = \bm{X}_{c}^\top \bm{y}_{c} \in \mathbb{R}^{p}, \quad
  s_{yy} = \bm{y}_{c}^\top \bm{y}_{c} \in \mathbb{R}.
\]

Let $\bm{A} = \bm{S}_{xx} + \lambda \bm{I}_p$.
Integrating out $(b_g, \bm{w}_g)$ yields a marginal likelihood (up to an additive constant):
\[
  \log p(\bm{y}_g \mid \bm{X}_g, \sigma_y^2, \lambda)
  = -\frac{1}{2}\left[
    (n_g - 1)\log \sigma_y^2
    + \log|\bm{A}|
    - p\log \lambda
    + \frac{1}{\sigma_y^2}\left(s_{yy} - \bm{s}_{xy}^\top \bm{A}^{-1}\bm{s}_{xy}\right)
  \right].
\]

This is the quantity implemented in PyMC as a \texttt{pm.Potential}, with all terms computed from the per-group sufficient statistics.

\subsection{Middle ground for readability: explicit intercept, collapsed slopes}

The fully collapsed formulation integrates out both $(b_g, \bm{w}_g)$, which is fast but can be hard to follow because the intercept
does not appear as an explicit node in the PyMC model.

As a more readable middle ground, we can keep the per-group intercepts $b_g$ as explicit PyMC variables, while still integrating out
the (high-dimensional) slopes $\bm{w}_g$ analytically. This preserves most of the scalability benefits (we still avoid sampling a
large $\bm{w}$ tensor), but makes it easier to extend the model with intercept pooling later (e.g., cluster-level intercepts).

In code this corresponds to \texttt{--intercept-mode explicit} in
\texttt{scripts/pipelines/train\_rt\_pymc\_collapsed\_ridge.py}.

For a group $g$ with uncentered design $\bm{X}_g \in \mathbb{R}^{n_g \times p}$ and response $\bm{y}_g \in \mathbb{R}^{n_g}$,
define $\bm{A} = \bm{X}_g^\top \bm{X}_g + \lambda \bm{I}_p$ and residual $\bm{r}_g = \bm{y}_g - b_g\bm{1}$.
Integrating out $\bm{w}_g$ yields:
\[
  \log p(\bm{y}_g \mid \bm{X}_g, b_g, \sigma_y^2, \lambda)
  = -\frac{1}{2}\left[
    n_g \log \sigma_y^2
    + \log|\bm{A}|
    - p\log \lambda
    + \frac{1}{\sigma_y^2}\left(\bm{r}_g^\top \bm{r}_g - (\bm{X}_g^\top \bm{r}_g)^\top \bm{A}^{-1}(\bm{X}_g^\top \bm{r}_g)\right)
  \right].
\]

All required terms can still be computed from per-group sufficient statistics:
\[
  \bm{r}_g^\top \bm{r}_g = \bm{y}_g^\top \bm{y}_g - 2 b_g (\bm{1}^\top \bm{y}_g) + n_g b_g^2,\quad
  \bm{X}_g^\top \bm{r}_g = \bm{X}_g^\top \bm{y}_g - b_g (\bm{X}_g^\top \bm{1}).
\]

\subsection{Per-group posterior coefficients and predictive intervals}

Conditioned on $(\sigma_y^2, \lambda)$, the posterior for slopes is:
\[
  \bm{w}_g \mid \bm{y}_g, \bm{X}_g, \sigma_y^2, \lambda
  \sim \mathcal{N}\!\left(\hat{\bm{w}}_g, \sigma_y^2 \bm{A}^{-1}\right),
  \qquad
  \hat{\bm{w}}_g = \bm{A}^{-1}\bm{s}_{xy}.
\]

The intercept is a derived random variable:
\[
  b_g = \bar{y}_g - \bar{\bm{x}}_g^\top \bm{w}_g,
\]
so:
\[
  \mathbb{E}[b_g] = \bar{y}_g - \bar{\bm{x}}_g^\top \hat{\bm{w}}_g,
  \quad
  \mathrm{Var}(b_g) = \sigma_y^2\left(\frac{1}{n_g} + \bar{\bm{x}}_g^\top \bm{A}^{-1}\bar{\bm{x}}_g\right),
  \quad
  \mathrm{Cov}(b_g, \bm{w}_g) = -\sigma_y^2 \bm{A}^{-1}\bar{\bm{x}}_g.
\]

We write these as a full coefficient covariance $\widehat{\Sigma}_{\beta,g}$ for
$\bm{\beta}_g = [b_g; \bm{w}_g]$, stored in \texttt{Stage1CoeffSummaries} as \texttt{beta\_cov}.

For a new row with features $\bm{x}$ (after the stored global transform), define $\bm{x}_1 = [1; \bm{x}]$.
The posterior predictive distribution is approximated as Normal:
\[
  y^\ast \mid \bm{x}, \text{data} \approx \mathcal{N}\!\left(
    \bm{x}_1^\top \hat{\bm{\beta}}_g,\;
    \sigma_y^2 + \bm{x}_1^\top \widehat{\Sigma}_{\beta,g}\bm{x}_1
  \right).
\]

\subsection{Planned extension: shared chemistry intercept hierarchy (partial collapse)}

\paragraph{Problem.}
The current baseline is evaluated on \emph{seen groups} only: we require that the group key
$g = (\texttt{species\_cluster}, \texttt{comp\_id})$ exists in the trained artifact, and we skip any rows whose group is missing.
This is appropriate for benchmarking against per-group ridge, but it does not answer a common production need:
\textbf{what do we predict for a new compound (or new group) that is not present in training?}

\paragraph{Motivation.}
We want a model that:
\begin{itemize}
  \item stays competitive with the current ridge-equivalent baseline on \emph{seen} $(\texttt{species\_cluster}, \texttt{comp\_id})$ groups;
  \item supports a principled backoff for \emph{new compounds}, ideally leveraging chemistry similarity so ``similar compounds'' can share strength;
  \item produces usable predictive intervals for downstream peak assignment; and
  \item is easier to read and extend than a fully-collapsed derivation (i.e., make the hierarchy explicit in the PyMC graph while
        still keeping training scalable).
\end{itemize}

\paragraph{Design principle.}
Do \emph{not} pool indiscriminately across \texttt{species\_cluster} (clusters correspond to substantially different matrices and/or
experimental conditions), but \emph{do} allow chemistry information to be shared globally with strong shrinkage so that a chemical baseline
can generalize to new compounds.

\paragraph{Plan.}
We keep the computationally critical part of the baseline (collapsing the ridge slopes) but make intercept structure explicit:
\begin{itemize}
  \item Keep per-group intercepts $b_g$ as explicit PyMC variables (as in \texttt{--intercept-mode explicit}).
  \item Add a hierarchical prior on $b_g$ that decomposes it into:
        (i) a cluster-specific offset, and (ii) a shared chemistry baseline.
  \item Continue to integrate out the per-group slopes $\bm{w}_g$ analytically via a collapsed Gaussian likelihood (so we avoid sampling
        a large $G \times p$ coefficient tensor).
\end{itemize}

\subsubsection{Revised model specification}

\paragraph{Indices and data.}
Let:
\begin{itemize}
  \item $c \in \{1,\dots,C\}$ index \texttt{species\_cluster} (matrix/condition clusters);
  \item $k \in \{1,\dots,K\}$ index chemicals (\texttt{chem\_id} = \texttt{compound} in the production CSV);
  \item $g \in \{1,\dots,G\}$ index groups defined by $g = (\texttt{species\_cluster}, \texttt{comp\_id})$;
  \item $k(g)$ map each group $g$ to its chemical id (resolved from \texttt{comp\_id} $\rightarrow$ \texttt{chem\_id} mapping);
  \item $c(g)$ map each group $g$ to its cluster id;
  \item $\bm{x}_{gi} \in \mathbb{R}^p$ be run covariates (IS*/RS*/ES\_* after the stored global transform);
  \item $y_{gi}$ be \texttt{rt} for row $i$ in group $g$.
\end{itemize}

Optionally, each chemical $k$ has a fixed feature/embedding vector $\bm{z}_k \in \mathbb{R}^d$
(e.g., a ChemBERTa embedding or other chemistry descriptor).

\paragraph{Chemistry baseline (deterministic regression).}
To keep inference tractable while preserving chemistry-driven generalization, we use a deterministic chemistry baseline:
\[
  t_0 \sim \mathcal{N}(0, s_0^2),
\]
\[
  t_k := t_0 + \bm{z}_k^\top \bm{\theta},
\]
and place a Normal prior on regression weights:
\[
  \bm{\theta} \sim \mathcal{N}(\bm{0}, s_{\theta}^2 \bm{I}_d).
\]
We use $\tau_t$ only as an additional uncertainty term for genuinely new chemicals at evaluation time, rather than sampling a separate
random effect $t_k$ for every chemical.

\paragraph{Cluster offsets (no cross-cluster pooling beyond a shared scale).}
Each cluster has its own additive offset:
\[
  \mu_c \mid \tau_\mu^2 \sim \mathcal{N}(0, \tau_\mu^2).
\]

\paragraph{Group intercepts (the ``seen-group head'').}
Each seen group $g$ has an intercept $b_g$ that shrinks toward the cluster offset plus the chemistry baseline:
\[
  b_g \mid \mu_{c(g)}, t_{k(g)}, \tau_b^2 \sim \mathcal{N}\!\left(\mu_{c(g)} + t_{k(g)}, \tau_b^2\right).
\]
Intuitively: $t_k$ captures shared chemistry structure; $\mu_c$ captures cluster-wide RT shifts; $\tau_b$ governs how much we allow
group-specific deviations around those shared terms.

\paragraph{Run-covariate slopes (ridge; collapsed).}
As in the baseline, each group has a slope vector:
\[
  \bm{w}_g \mid \sigma_y^2, \lambda \sim \mathcal{N}\!\left(\bm{0}, \frac{\sigma_y^2}{\lambda}\bm{I}_p\right),
\]
and observations follow:
\[
  y_{gi} = b_g + \bm{x}_{gi}^\top \bm{w}_g + \varepsilon_{gi}, \qquad
  \varepsilon_{gi} \sim \mathcal{N}(0, \sigma_y^2).
\]

\paragraph{Hyperpriors.}
We will use weakly-informative priors on scale parameters in RT units:
\[
  \sigma_y \sim \mathrm{HalfNormal}(s_{\sigma}),\quad
  \tau_\mu \sim \mathrm{HalfNormal}(s_\mu),\quad
  \tau_b \sim \mathrm{HalfNormal}(s_b),
\]
and a broad Normal prior on the global baseline:
\[
  t_0 \sim \mathcal{N}(0, s_0^2),
\]
and, if chemistry features are used, a Normal prior on regression weights (as above).
The ridge penalty $\lambda$ can remain fixed (ridge-equivalent) or be learned as in \texttt{--lambda-mode learn}.

\subsubsection{How we keep training scalable (partial collapse via sufficient statistics)}

The PyMC graph will include $(\sigma_y, \lambda, \{\mu_c\}, \{t_k\}, \{b_g\}, \bm{\theta}, \tau_\cdot)$ explicitly.
We still integrate out $\bm{w}_g$ analytically, so the likelihood contribution per group remains a collapsed Gaussian term.

To evaluate the collapsed likelihood without row-level data, we use per-group sufficient statistics. The current trainer already computes
centered statistics from per-group means $(\bar{\bm{x}}_g, \bar{y}_g)$:
\[
  \bm{S}_{xx} = \bm{X}_c^\top \bm{X}_c,\quad
  \bm{s}_{xy} = \bm{X}_c^\top \bm{y}_c,\quad
  s_{yy} = \bm{y}_c^\top \bm{y}_c,\quad
  n_g.
\]
From these we can recover the uncentered statistics needed for an explicit intercept $b_g$:
\[
  \bm{X}_g^\top \bm{X}_g = \bm{S}_{xx} + n_g \bar{\bm{x}}_g \bar{\bm{x}}_g^\top,\quad
  \bm{X}_g^\top \bm{y}_g = \bm{s}_{xy} + \bar{\bm{x}}_g (n_g \bar{y}_g),
\]
\[
  \bm{X}_g^\top \bm{1} = n_g \bar{\bm{x}}_g,\quad
  \bm{1}^\top \bm{y}_g = n_g \bar{y}_g,\quad
  \bm{y}_g^\top \bm{y}_g = s_{yy} + n_g \bar{y}_g^2.
\]

Given $b_g$, define the residual vector $\bm{r}_g = \bm{y}_g - b_g\bm{1}$ and let
\[
  \bm{A}_g = \bm{X}_g^\top \bm{X}_g + \lambda \bm{I}_p.
\]
Integrating out $\bm{w}_g$ gives the per-group collapsed log-likelihood (up to a constant):
\[
  \log p(\bm{y}_g \mid \bm{X}_g, b_g, \sigma_y^2, \lambda)
  = -\frac{1}{2}\left[
    n_g \log \sigma_y^2
    + \log|\bm{A}_g|
    - p\log \lambda
    + \frac{1}{\sigma_y^2}\left(\bm{r}_g^\top \bm{r}_g - (\bm{X}_g^\top \bm{r}_g)^\top \bm{A}_g^{-1}(\bm{X}_g^\top \bm{r}_g)\right)
  \right].
\]
This is the only ``linear algebra heavy'' component, and it remains unchanged from the baseline except that $b_g$ is now explicit.

\subsubsection{Inference and production artifacts}

\paragraph{Inference strategy.}
This model introduces many explicit latent variables ($b_g$ for each group, $t_k$ for each chemical, and $\mu_c$ for each cluster).
To keep training practical, we plan to use the same inference modes already supported by the trainer:
\begin{itemize}
  \item MAP (fast default) for point estimates of hyperparameters and latent intercept terms;
  \item ADVI when we need approximate posterior draws to propagate uncertainty more fully; and
  \item NUTS only for debugging on small subsets (e.g., limited number of groups) to sanity-check geometry and priors.
\end{itemize}
In all cases, the slopes $\bm{w}_g$ remain collapsed analytically.

\paragraph{What must be saved for production use.}
For \emph{seen groups}, we can continue to materialize a \texttt{Stage1CoeffSummaries} artifact (per-group $\hat{\bm{\beta}}_g$ and
covariance) so the existing streaming evaluator can produce predictions and intervals.

To support \emph{backoff predictions for new compounds}, the production artifact must additionally include posterior summaries for:
\begin{itemize}
  \item cluster offsets $\{\mu_c\}$ (and their uncertainty),
  \item chemical baselines $\{t_k\}$ (and their uncertainty),
  \item the global baseline $t_0$,
  \item chemistry regression weights $\bm{\theta}$ (and their uncertainty) if $\bm{z}_k$ features are used,
  \item the scale terms $(\tau_t, \tau_\mu, \tau_b)$ that govern cold-start uncertainty, and
  \item the global noise and ridge hyperparameters $(\sigma_y, \lambda)$.
\end{itemize}
This can be implemented either by extending the existing \texttt{Stage1CoeffSummaries} format or by writing a companion artifact
containing these additional parameters. The evaluator will need a corresponding ``backoff'' path for rows whose group key is missing.

\subsubsection{Predictions and intervals (seen groups and backoff for new compounds)}

\paragraph{Seen group prediction.}
For a seen group $g$, the predictive mean for a row with covariates $\bm{x}$ is:
\[
  \mathbb{E}[y^\ast \mid \bm{x}, \text{posterior}] \approx \mathbb{E}\left[b_g + \bm{x}^\top \hat{\bm{w}}_g\right],
\]
where for fixed $(b_g, \sigma_y^2, \lambda)$ the ridge posterior mean is:
\[
  \hat{\bm{w}}_g = \bm{A}_g^{-1}\bm{X}_g^\top(\bm{y}_g - b_g\bm{1}).
\]
We will propagate uncertainty by drawing from the PyMC posterior over the explicit variables and, for each draw, using the analytic
Gaussian form of $\bm{w}_g \mid b_g, \sigma_y^2, \lambda$.

\paragraph{Predictive intervals.}
Conditional on $(b_g, \sigma_y^2, \lambda)$, the slopes have covariance $\sigma_y^2 \bm{A}_g^{-1}$ and the predictive variance for a
new row is:
\[
  \mathrm{Var}(y^\ast \mid \bm{x}, b_g, \sigma_y^2, \lambda)
  \approx \sigma_y^2 + \bm{x}^\top (\sigma_y^2 \bm{A}_g^{-1}) \bm{x}.
\]
Uncertainty in $(b_g, \mu_c, t_k, \bm{\theta}, \tau_\cdot, \lambda, \sigma_y)$ is then incorporated by Monte Carlo over posterior draws,
yielding calibrated intervals for downstream usage.

\paragraph{Backoff for new compounds (new $\texttt{comp\_id}$ and/or new chemical id).}
For a new compound not present in training, we do not have a learned $b_g$ for its $(c, \texttt{comp\_id})$ key. The hierarchy provides
a principled backoff:
\[
  b_{\text{new}} \sim \mathcal{N}\!\left(\mu_c + t_k, \tau_b^2\right),
\]
where $t_k$ is either:
\begin{itemize}
  \item drawn from the learned chemical baseline if the chemical id $k$ was seen in training, or
  \item predicted from chemistry features via $t_k = t_0 + \bm{z}_k^\top \bm{\theta}$ for a genuinely new chemical, with additional
        uncertainty $\tau_t^2$.
\end{itemize}
This supports cold-start predictions that share strength across chemically similar compounds while respecting matrix differences via $\mu_c$.
In the base formulation above, unseen-group slopes $\bm{w}$ revert to the ridge prior mean $\bm{0}$; if run-covariate response is required
for cold-start accuracy, we can extend the model later with structured slope means (e.g., cluster-level or chemistry-gated slope means)
while keeping the same partial-collapsed strategy.

\section{Benchmark results (cap100 train $\rightarrow$ realtest eval)}

All numbers below train on the full cap100 CSV and evaluate on the corresponding realtest CSV used by
\texttt{scripts/run\_rt\_prod\_eval.sh}. Evaluation is restricted to rows whose group key
$(\texttt{species\_cluster}, \texttt{comp\_id})$ exists in the coefficient artifact (seen groups only).

\begin{table}[H]
\centering
\begin{tabular}{llrrr}
\toprule
Lib & Model & $n$ used & RMSE & cov$_{0.95}$ \\
\midrule
208 & Lasso (eslasso) & 915{,}327 & 0.0151 & 0.954 \\
208 & PyMC chem\_hier (ADVI; explicit intercepts; $\lambda$ fixed) & 966{,}482 & 0.0094 & 0.984 \\
\midrule
209 & Lasso (eslasso) & 13{,}301{,}888 & 0.0094 & 0.993 \\
209 & PyMC chem\_hier (ADVI; explicit intercepts; $\lambda$ fixed) & 13{,}596{,}681 & 0.0085 & 0.990 \\
\bottomrule
\end{tabular}
\caption{Global metrics (trained on full cap100; evaluated on realtest; seen groups only).}
\end{table}

\paragraph{Runtime (end-to-end).}
Training the chem\_hier variant (explicit intercepts; ADVI with 5000 steps) took approximately
$27$m for lib208 and $48$m for lib209 on the current workstation. Realtest evaluation took approximately $17$s/$214$s using posterior
predictive intervals (200 draws).

\paragraph{Predictive intervals (cov$_{0.95}$).}
The reported cov$_{0.95}$ is the fraction of test rows whose true RT falls inside the nominal 95\% predictive interval.
For PyMC chem\_hier, we compute this via posterior predictive Monte Carlo
(see \texttt{scripts/pipelines/eval\_rt\_pymc\_posterior\_predictive.py}).
For the lasso baseline, the ``interval'' is a model-provided window; cov$_{0.95}$ is the fraction within that window.

\subsection{Partially-collapsed chem\_hier vs lasso baseline (realtest)}

This comparison mirrors \texttt{scripts/run\_rt\_prod\_eval.sh}: we aggregate RMSE by
\texttt{species\_group\_raw} and show compound coverage (unique \texttt{compound} ids) alongside RMSE.

\begin{table}[H]
\centering
\begin{tabular}{p{0.44\textwidth}rrrr}
\toprule
Species group & RMSE (chem\_hier) & RMSE (lasso) & Ratio & Lasso cov. \\
\midrule
1 human blood & 0.0094 & 0.0123 & 0.77 & 0.90 \\
2 mammal / rodent blood & 0.0136 & 0.0215 & 0.63 & 0.89 \\
3 human / mammal fecal & 0.0075 & 0.0217 & 0.35 & 0.86 \\
4 human / animal urine & 0.0100 & 0.0156 & 0.64 & 0.82 \\
5 mammal tissue & 0.0083 & 0.0129 & 0.64 & 0.78 \\
6 non-human cells + plants & 0.0078 & 0.0127 & 0.61 & 0.79 \\
8 human cells & 0.0065 & 0.0196 & 0.33 & 0.76 \\
\bottomrule
\end{tabular}
\caption{lib208 realtest: RMSE by \texttt{species\_group\_raw} (PyMC chem\_hier vs lasso).}
\end{table}

\begin{table}[H]
\centering
\begin{tabular}{p{0.44\textwidth}rrrr}
\toprule
Species group & RMSE (chem\_hier) & RMSE (lasso) & Ratio & Lasso cov. \\
\midrule
1 human blood & 0.0076 & 0.0078 & 0.98 & 0.94 \\
2 mammal / rodent blood & 0.0077 & 0.0101 & 0.76 & 0.91 \\
3 human / mammal fecal & 0.0114 & 0.0124 & 0.92 & 0.94 \\
4 human / fecal urine & 0.0093 & 0.0115 & 0.81 & 0.91 \\
5 mammal tissue & 0.0101 & 0.0120 & 0.84 & 0.81 \\
6 non-human cells + plants & 0.0106 & 0.0147 & 0.72 & 0.80 \\
8 human cells & 0.0065 & 0.0126 & 0.52 & 0.83 \\
\bottomrule
\end{tabular}
\caption{lib209 realtest: RMSE by \texttt{species\_group\_raw} (PyMC chem\_hier vs lasso).}
\end{table}

\begin{figure}[H]
  \centering
  \begin{subfigure}[b]{0.95\textwidth}
    \centering
    \includegraphics[width=\textwidth]{images/rt_pymc_collapsed_ridge_report/rt_compare_rmse_by_species_group_lib208_realtest_rmse_coverage_by_species_group.png}
    \caption{lib208 (realtest).}
  \end{subfigure}
  \vspace{0.5em}
  \begin{subfigure}[b]{0.95\textwidth}
    \centering
    \includegraphics[width=\textwidth]{images/rt_pymc_collapsed_ridge_report/rt_compare_rmse_by_species_group_lib209_realtest_rmse_coverage_by_species_group.png}
    \caption{lib209 (realtest).}
  \end{subfigure}
  \caption{RMSE (top) and compound coverage (bottom) by \texttt{species\_group\_raw}: PyMC chem\_hier vs lasso baseline.}
\end{figure}

\subsection{Cap100 coverage (per species$\times$compound) and implications for pooling}

\paragraph{Cap semantics.}
The capped datasets under \texttt{repo\_export/lib*/cap*} are built by down-sampling rows per
$(\texttt{species\_matrix\_type}, \texttt{species}, \texttt{comp\_id})$ key (see
\texttt{scripts/pipelines/sample\_per\_species\_compound.py}).
Therefore, ``cap100'' does \emph{not} mean each compound has 100 rows; it means each
species$\times$compound pair has \emph{at most} 100 rows, and a compound's total rows can be far above (or below) 100 depending on
its species coverage and raw availability.

\paragraph{Empirical cap100 coverage (lib208/lib209).}
Using the Parquets under \texttt{repo\_export/lib{208,209}/cap100} with the corresponding
\texttt{repo\_export/lib*/species\_mapping/*\_species\_mapping.csv}, we computed per-key coverage:
\begin{itemize}
  \item \textbf{lib208 cap100}: 1,379 compounds, 25 species, 16,196 species$\times$compound pairs; pair median = 96 rows, p90 = 100;
        49.0\% of pairs hit cap exactly (51.0\% are under-filled). Median species-per-compound = 6 (p90 = 20).
  \item \textbf{lib209 cap100}: 2,640 compounds, 103 species, 51,943 species$\times$compound pairs; pair median = 68 rows, p90 = 100;
        40.6\% of pairs hit cap exactly (59.4\% are under-filled). Median species-per-compound = 7 (p90 = 63).
\end{itemize}
At the compound level (summing across species), cap100 still yields large totals on average (median $\approx$ 455--536 rows per
compound), but there is a consistent long tail: roughly 12--14\% of compounds have $<100$ total rows.

\paragraph{Why this matters for model design (seen compounds).}
For the ``seen chemical'' setting (predict new runs for compounds present in training), coverage is generally strong for many
compound$\times$species pairs, but there is a sizeable long tail of sparse pairs and sparse compounds.
This suggests that the most useful pooling is \emph{not} chemistry-based (cold start) but rather shrinkage that stabilizes estimates
for under-observed groups while leaving well-observed groups effectively unpooled:
\begin{itemize}
  \item \textbf{Random slopes with hierarchical shrinkage:} per-compound (or per group) run-covariate coefficients shrunk toward a
        global mean, improving stability for sparse compounds/pairs without degrading well-covered ones.
  \item \textbf{Optional pair-level intercept pooling:} an additional intercept term for $(\texttt{species\_cluster},\texttt{comp\_id})$
        (or species$\times$compound) with shrinkage toward the compound intercept, to stabilize sparse species$\times$compound pairs.
\end{itemize}
We intentionally avoid run/batch random intercepts for this report because they do not extrapolate to new runs unless the run id is
known at prediction time; instead, generalization to new runs is driven by observed run covariates (\texttt{IS*}/\texttt{RS*}/\texttt{ES\_*}).

\subsection{Seen-group ablations: do we need chemistry embeddings or pooling?}

\paragraph{Question.}
For the ``seen group'' setting (predict new runs for $(\texttt{species\_cluster}, \texttt{comp\_id})$ groups present in training),
do we gain anything from:
\begin{itemize}
  \item adding chemistry embeddings (ChemBERTa PCA-20) to the intercept hierarchy, or
  \item adding additional pooling structures (per-compound intercept pooling, per-compound slope pooling)?
\end{itemize}

\paragraph{Result (cap100 $\rightarrow$ realtest; seen groups only).}
On realtest, chemistry embeddings and the additional pooling variants below do \emph{not} improve global RMSE relative to the simplest
collapsed ridge baseline. The largest RMSE reduction comes from adding a small nonlinearity to the anchor predictors
(\texttt{IS*}/\texttt{RS*} poly2 expansion).

For uncertainty, the per-group coefficient-summary baseline under-covers when using the analytic Normal approximation directly.
Two simple fixes yield near-nominal coverage:
(a) post-hoc scaling of predictive standard deviation (\texttt{--pred-std-scale}), or
(b) using the PyMC collapsed ridge trainer, which estimates a global $\sigma_y$ and tends to give more conservative intervals.

\begin{table}[H]
\centering
\begin{tabular}{llrr}
\toprule
Lib & Model (cap100 train) & RMSE & cov$_{0.95}$ \\
\midrule
208 & sklearn ridge + poly2(\texttt{IS*}/\texttt{RS*}) anchors (two-$\lambda$) & 0.00891 & 0.929 \\
208 & sklearn ridge + poly2 anchors + pred\_std\_scale=1.135 & 0.00891 & 0.951 \\
208 & PyMC collapsed ridge + poly2 anchors (MAP; fixed two-$\lambda$) & 0.00891 & 0.945 \\
208 & collapsed ridge (MAP; flat intercept) & 0.00934 & 0.946 \\
208 & explicit intercept + chem\_hier (ADVI; ChemBERTa PCA-20) & 0.00934 & 0.984 \\
208 & explicit intercept + comp\_hier (MAP; no chemistry) & 0.00938 & 0.947 \\
208 & comp\_hier + comp\_id slope head (MAP) & 0.00938 & 0.946 \\
\midrule
209 & sklearn ridge + poly2(\texttt{IS*}/\texttt{RS*}) anchors (two-$\lambda$) & 0.00800 & 0.913 \\
209 & sklearn ridge + poly2 anchors + pred\_std\_scale=1.198 & 0.00800 & 0.950 \\
209 & PyMC collapsed ridge + poly2 anchors (MAP; fixed two-$\lambda$) & 0.00800 & 0.958 \\
209 & collapsed ridge (MAP; flat intercept) & 0.00833 & 0.963 \\
209 & explicit intercept + chem\_hier (ADVI; ChemBERTa PCA-20) & 0.00838 & 0.991 \\
209 & explicit intercept + comp\_hier (ADVI 3000) & 0.00853 & 0.998 \\
209 & comp\_hier + comp\_id slope head (ADVI 3000) & 0.00845 & 0.998 \\
\bottomrule
\end{tabular}
\caption{Seen-group ablations on realtest (cap100 train $\rightarrow$ realtest eval; seen groups only).}
\end{table}

\paragraph{Takeaway.}
At cap100, most groups have enough observations that per-group fits already generalize well to realtest, so additional partial pooling
does not move global RMSE. If we want to keep the PyMC model minimal for the seen-group project, the best-performing choice is the
collapsed ridge baseline (no chemistry embeddings). However, a simple nonlinear anchor expansion achieves a meaningful RMSE reduction on
both libraries, and does \emph{not} require PyMC.

If calibrated uncertainty is a primary goal, the per-group closed-form baseline needs a variance calibration step. In practice this can
be done with a single post-hoc scale factor on predictive standard deviation (estimated on a validation set), or by using the PyMC
collapsed ridge trainer with fixed $\lambda$ which gives similar point predictions while typically producing closer-to-nominal coverage
via a global $\sigma_y$ estimate.

We also tried learning $\lambda$ via PyMC (MAP over the collapsed marginal likelihood). In this setting it tended to over-shrink the
poly2 terms (very large $\lambda_{\text{poly}}$) and regress toward the linear ridge baseline, worsening RMSE. Therefore, for RMSE we
prefer fixed $\lambda$ tuned for prediction (e.g., by a small grid search), and use PyMC primarily for coherent uncertainty or future
model extensions.

\subsection{Held-out chemicals (cap100 train; cold-start backoff on unseen groups)}

\paragraph{Goal.}
Quantify what the chem\_hier intercept hierarchy buys us for \emph{unseen chemicals} (cold start), since realtest contains only seen
$(\texttt{species\_cluster}, \texttt{comp\_id})$ groups for the cap100 training artifacts.

\paragraph{Experiment.}
For each library independently, we:
\begin{itemize}
  \item randomly hold out 10\% of unique \texttt{compound} (chem\_id) values from the cap100 training CSV (seed 42);
  \item train chem\_hier on the remaining rows (ADVI 5000 steps; $\lambda$ fixed);
  \item evaluate on the held-out slice without requiring seen groups, so every row is scored via the backoff model
        (unseen $(\texttt{species\_cluster}, \texttt{comp\_id})$ by construction).
\end{itemize}

\paragraph{Backoff variants compared.}
\begin{itemize}
  \item \textbf{chem\_hier backoff (chem)}: $b = \mu_{c} + t_0 + \bm{z}_k^\top\bm{\theta}$ for unseen chemicals.
  \item \textbf{cluster-only backoff}: $b = \mu_{c} + t_0$ (ignore chemistry; same uncertainty terms).
\end{itemize}
In both cases, unseen-group slopes revert to the ridge prior mean $\bm{0}$ (no shared slope head); this is a strict zero-shot setting.

\begin{table}[H]
\centering
\begin{tabular}{llrrr}
\toprule
Lib & Backoff mean & $n$ used & RMSE & cov$_{0.95}$ \\
\midrule
208 & chem\_hier (chem) & 121{,}092 & 1.3959 & 1.000 \\
208 & cluster-only & 121{,}092 & 1.4078 & 0.999 \\
\midrule
209 & chem\_hier (chem) & 357{,}488 & 1.8529 & 1.000 \\
209 & cluster-only & 357{,}488 & 1.9604 & 0.999 \\
\bottomrule
\end{tabular}
\caption{Held-out chemicals (10\% chem\_id holdout from cap100; backoff-only evaluation on the held-out slice).}
\end{table}

\paragraph{RMSE by species\_group\_raw (held-out chemicals).}
\begin{figure}[H]
  \centering
  \begin{subfigure}[b]{0.95\textwidth}
    \centering
    \includegraphics[width=\textwidth]{images/rt_pymc_collapsed_ridge_report/rt_compare_rmse_by_species_group_lib208_holdout_chems_rmse_coverage_by_species_group.png}
    \caption{lib208 (held-out chemicals).}
  \end{subfigure}
  \vspace{0.5em}
  \begin{subfigure}[b]{0.95\textwidth}
    \centering
    \includegraphics[width=\textwidth]{images/rt_pymc_collapsed_ridge_report/rt_compare_rmse_by_species_group_lib209_holdout_chems_rmse_coverage_by_species_group.png}
    \caption{lib209 (held-out chemicals).}
  \end{subfigure}
  \caption{Held-out chemicals: chem\_hier backoff vs cluster-only backoff by \texttt{species\_group\_raw}.}
\end{figure}

\section{Reproduction commands (cap100 train $\rightarrow$ realtest eval)}

\subsection{Per-species-group comparison (chem\_hier vs lasso; realtest)}
\begin{verbatim}
# Train chem_hier (cap100)
poetry run python scripts/pipelines/train_rt_pymc_collapsed_ridge.py \
  --data-csv repo_export/lib208/cap100/merged_training_all_lib208_cap100_chemclass_rt_prod.csv \
  --include-es-all --lambda-mode fixed --lambda-slopes 1e-3 \
  --intercept-mode explicit --intercept-prior chem_hier \
  --chem-embeddings-path resources/metabolites/embeddings_chemberta_pca20.parquet \
  --chem-feature-center global \
  --method advi --advi-steps 5000 --advi-draws 200 \
  --output-dir output/tmp_pymc_chem_hier_lib208_cap100_advi --seed 42

poetry run python scripts/pipelines/train_rt_pymc_collapsed_ridge.py \
  --data-csv repo_export/lib209/cap100/merged_training_all_lib209_cap100_chemclass_rt_prod.csv \
  --include-es-all --lambda-mode fixed --lambda-slopes 1e-3 \
  --intercept-mode explicit --intercept-prior chem_hier \
  --chem-embeddings-path resources/metabolites/embeddings_chemberta_pca20.parquet \
  --chem-feature-center global \
  --method advi --advi-steps 5000 --advi-draws 200 \
  --output-dir output/tmp_pymc_chem_hier_lib209_cap100_advi --seed 42

# Lasso baseline aggregated by species_group_raw
poetry run python scripts/pipelines/eval_rt_lasso_baseline.py \
  --train-output-dir output/tmp_lasso_lib208_realtest --lib-id 208 \
  --test-csv repo_export/lib208/realtest/merged_training_realtest_lib208_chemclass_rt_prod.csv \
  --label realtest

poetry run python scripts/pipelines/eval_rt_lasso_baseline.py \
  --train-output-dir output/tmp_lasso_lib209_realtest --lib-id 209 \
  --test-csv repo_export/lib209/realtest/merged_training_realtest_lib209_chemclass_rt_prod.csv \
  --label realtest

# chem_hier aggregated by species_group_raw
poetry run python scripts/pipelines/eval_rt_pymc_posterior_predictive.py \
  --train-output-dir output/tmp_pymc_chem_hier_lib208_cap100_advi \
  --test-csv repo_export/lib208/realtest/merged_training_realtest_lib208_chemclass_rt_prod.csv \
  --require-seen-group --write-by-species-group \
  --max-test-rows 0 --chunk-size 50000 --n-ppc-draws 200 \
  --label realtest_seen

poetry run python scripts/pipelines/eval_rt_pymc_posterior_predictive.py \
  --train-output-dir output/tmp_pymc_chem_hier_lib209_cap100_advi \
  --test-csv repo_export/lib209/realtest/merged_training_realtest_lib209_chemclass_rt_prod.csv \
  --require-seen-group --write-by-species-group \
  --max-test-rows 0 --chunk-size 50000 --n-ppc-draws 200 \
  --label realtest_seen

# Side-by-side plot by species_group_raw
poetry run python scripts/pipelines/compare_rt_models_by_group.py \
  --hier-csv output/tmp_pymc_chem_hier_lib208_cap100_advi/results/rt_eval_ppc_by_species_group_realtest_seen.csv \
  --lasso-csv output/tmp_lasso_lib208_realtest/results/rt_eval_lasso_by_species_group_realtest.csv \
  --hier-label "PyMC chem_hier" --lasso-label "Lasso" --label realtest \
  --output docs/models/images/rt_pymc_collapsed_ridge_report/rt_compare_rmse_by_species_group_lib208_realtest.png

poetry run python scripts/pipelines/compare_rt_models_by_group.py \
  --hier-csv output/tmp_pymc_chem_hier_lib209_cap100_advi/results/rt_eval_ppc_by_species_group_realtest_seen.csv \
  --lasso-csv output/tmp_lasso_lib209_realtest/results/rt_eval_lasso_by_species_group_realtest.csv \
  --hier-label "PyMC chem_hier" --lasso-label "Lasso" --label realtest \
  --output docs/models/images/rt_pymc_collapsed_ridge_report/rt_compare_rmse_by_species_group_lib209_realtest.png
\end{verbatim}

\subsection{Nonlinear anchor baseline (poly2 on \texttt{IS*}/\texttt{RS*}; realtest)}
\begin{verbatim}
# lib208: poly2 anchor expansion (cap100 train -> realtest eval)
poetry run python scripts/pipelines/train_rt_stage1_coeff_summaries.py \
  --data-csv repo_export/lib208/cap100/merged_training_all_lib208_cap100_chemclass_rt_prod.csv \
  --include-es-all --feature-center none --feature-rotation none \
  --lambda-ridge 3e-4 --anchor-expansion poly2 --lambda-ridge-poly 1e-4 \
  --output-dir output/tmp_stage1_ridge_poly2_lib208_cap100 --seed 42

poetry run python scripts/pipelines/eval_rt_coeff_summaries_by_species_cluster.py \
  --coeff-npz output/tmp_stage1_ridge_poly2_lib208_cap100/stage1_coeff_summaries.npz \
  --test-csv repo_export/lib208/realtest/merged_training_realtest_lib208_chemclass_rt_prod.csv \
  --chunk-size 50000 --max-test-rows 0 --label realtest

# Optional: suggest and apply a variance scale factor for nominal coverage
poetry run python scripts/pipelines/eval_rt_coeff_summaries_by_species_cluster.py \
  --coeff-npz output/tmp_stage1_ridge_poly2_lib208_cap100/stage1_coeff_summaries.npz \
  --test-csv repo_export/lib208/realtest/merged_training_realtest_lib208_chemclass_rt_prod.csv \
  --chunk-size 50000 --max-test-rows 0 --label realtest --suggest-pred-std-scale

poetry run python scripts/pipelines/eval_rt_coeff_summaries_by_species_cluster.py \
  --coeff-npz output/tmp_stage1_ridge_poly2_lib208_cap100/stage1_coeff_summaries.npz \
  --test-csv repo_export/lib208/realtest/merged_training_realtest_lib208_chemclass_rt_prod.csv \
  --chunk-size 50000 --max-test-rows 0 --label realtest_scaled --pred-std-scale 1.135

# lib209: poly2 anchor expansion (cap100 train -> realtest eval)
poetry run python scripts/pipelines/train_rt_stage1_coeff_summaries.py \
  --data-csv repo_export/lib209/cap100/merged_training_all_lib209_cap100_chemclass_rt_prod.csv \
  --include-es-all --feature-center none --feature-rotation none \
  --lambda-ridge 3e-4 --anchor-expansion poly2 --lambda-ridge-poly 1e-4 \
  --output-dir output/tmp_stage1_ridge_poly2_lib209_cap100 --seed 42

poetry run python scripts/pipelines/eval_rt_coeff_summaries_by_species_cluster.py \
  --coeff-npz output/tmp_stage1_ridge_poly2_lib209_cap100/stage1_coeff_summaries.npz \
  --test-csv repo_export/lib209/realtest/merged_training_realtest_lib209_chemclass_rt_prod.csv \
  --chunk-size 50000 --max-test-rows 0 --label realtest

poetry run python scripts/pipelines/eval_rt_coeff_summaries_by_species_cluster.py \
  --coeff-npz output/tmp_stage1_ridge_poly2_lib209_cap100/stage1_coeff_summaries.npz \
  --test-csv repo_export/lib209/realtest/merged_training_realtest_lib209_chemclass_rt_prod.csv \
  --chunk-size 50000 --max-test-rows 0 --label realtest --suggest-pred-std-scale

poetry run python scripts/pipelines/eval_rt_coeff_summaries_by_species_cluster.py \
  --coeff-npz output/tmp_stage1_ridge_poly2_lib209_cap100/stage1_coeff_summaries.npz \
  --test-csv repo_export/lib209/realtest/merged_training_realtest_lib209_chemclass_rt_prod.csv \
  --chunk-size 50000 --max-test-rows 0 --label realtest_scaled --pred-std-scale 1.198

# PyMC collapsed ridge (same poly2 nonlinearity; global sigma_y)
poetry run python scripts/pipelines/train_rt_pymc_collapsed_ridge.py \
  --data-csv repo_export/lib208/cap100/merged_training_all_lib208_cap100_chemclass_rt_prod.csv \
  --include-es-all --feature-center none --feature-rotation none \
  --anchor-expansion poly2 --lambda-mode fixed --lambda-slopes 3e-4 --lambda-slopes-poly 1e-4 \
  --intercept-mode collapsed --method map --output-dir output/tmp_rt_pymc_poly2_lib208_cap100_fixed --seed 42

poetry run python scripts/pipelines/eval_rt_coeff_summaries_by_species_cluster.py \
  --coeff-npz output/tmp_rt_pymc_poly2_lib208_cap100_fixed/models/stage1_coeff_summaries_posterior.npz \
  --test-csv repo_export/lib208/realtest/merged_training_realtest_lib208_chemclass_rt_prod.csv \
  --chunk-size 50000 --max-test-rows 0 --label realtest

poetry run python scripts/pipelines/train_rt_pymc_collapsed_ridge.py \
  --data-csv repo_export/lib209/cap100/merged_training_all_lib209_cap100_chemclass_rt_prod.csv \
  --include-es-all --feature-center none --feature-rotation none \
  --anchor-expansion poly2 --lambda-mode fixed --lambda-slopes 3e-4 --lambda-slopes-poly 1e-4 \
  --intercept-mode collapsed --method map --output-dir output/tmp_rt_pymc_poly2_lib209_cap100_fixed --seed 42

poetry run python scripts/pipelines/eval_rt_coeff_summaries_by_species_cluster.py \
  --coeff-npz output/tmp_rt_pymc_poly2_lib209_cap100_fixed/models/stage1_coeff_summaries_posterior.npz \
  --test-csv repo_export/lib209/realtest/merged_training_realtest_lib209_chemclass_rt_prod.csv \
  --chunk-size 50000 --max-test-rows 0 --label realtest
\end{verbatim}

\subsection{Held-out chemicals (10\% chem\_id holdout; cap100)}
\begin{verbatim}
# Build a cold-start split (hold out chem_ids) from the cap100 training CSV
poetry run python scripts/pipelines/split_rt_prod_csv_holdout_chems.py \
  --data-csv repo_export/lib208/cap100/merged_training_all_lib208_cap100_chemclass_rt_prod.csv \
  --output-dir output/rt_pymc_chem_hier_bench/holdout_chems/lib208_cap100_seed42_frac0.1 \
  --holdout-frac 0.1 --seed 42 --overwrite

poetry run python scripts/pipelines/split_rt_prod_csv_holdout_chems.py \
  --data-csv repo_export/lib209/cap100/merged_training_all_lib209_cap100_chemclass_rt_prod.csv \
  --output-dir output/rt_pymc_chem_hier_bench/holdout_chems/lib209_cap100_seed42_frac0.1 \
  --holdout-frac 0.1 --seed 42 --overwrite

# Train chem_hier on the reduced train split
poetry run python scripts/pipelines/train_rt_pymc_collapsed_ridge.py \
  --data-csv output/rt_pymc_chem_hier_bench/holdout_chems/lib208_cap100_seed42_frac0.1/train.csv \
  --include-es-all --lambda-mode fixed --lambda-slopes 1e-3 \
  --intercept-mode explicit --intercept-prior chem_hier \
  --chem-embeddings-path resources/metabolites/embeddings_chemberta_pca20.parquet \
  --chem-feature-center global \
  --method advi --advi-steps 5000 --advi-draws 200 \
  --output-dir output/tmp_pymc_chem_hier_lib208_cap100_holdoutchems_advi --seed 42

poetry run python scripts/pipelines/train_rt_pymc_collapsed_ridge.py \
  --data-csv output/rt_pymc_chem_hier_bench/holdout_chems/lib209_cap100_seed42_frac0.1/train.csv \
  --include-es-all --lambda-mode fixed --lambda-slopes 1e-3 \
  --intercept-mode explicit --intercept-prior chem_hier \
  --chem-embeddings-path resources/metabolites/embeddings_chemberta_pca20.parquet \
  --chem-feature-center global \
  --method advi --advi-steps 5000 --advi-draws 200 \
  --output-dir output/tmp_pymc_chem_hier_lib209_cap100_holdoutchems_advi --seed 42

# Evaluate on held-out chem_ids using backoff (unseen groups), and compare to a cluster-only backoff
poetry run python scripts/pipelines/eval_rt_pymc_posterior_predictive.py \
  --train-output-dir output/tmp_pymc_chem_hier_lib208_cap100_holdoutchems_advi \
  --backoff-npz output/tmp_pymc_chem_hier_lib208_cap100_holdoutchems_advi/models/chem_hier_backoff.npz \
  --test-csv output/rt_pymc_chem_hier_bench/holdout_chems/lib208_cap100_seed42_frac0.1/test.csv \
  --backoff-mean-mode chem \
  --max-test-rows 0 --chunk-size 50000 --n-ppc-draws 200 \
  --label holdout_chems --write-by-species-group

poetry run python scripts/pipelines/eval_rt_pymc_posterior_predictive.py \
  --train-output-dir output/tmp_pymc_chem_hier_lib208_cap100_holdoutchems_advi \
  --backoff-npz output/tmp_pymc_chem_hier_lib208_cap100_holdoutchems_advi/models/chem_hier_backoff.npz \
  --test-csv output/rt_pymc_chem_hier_bench/holdout_chems/lib208_cap100_seed42_frac0.1/test.csv \
  --backoff-mean-mode cluster_only \
  --max-test-rows 0 --chunk-size 50000 --n-ppc-draws 200 \
  --label holdout_chems_cluster_only --write-by-species-group

poetry run python scripts/pipelines/eval_rt_pymc_posterior_predictive.py \
  --train-output-dir output/tmp_pymc_chem_hier_lib209_cap100_holdoutchems_advi \
  --backoff-npz output/tmp_pymc_chem_hier_lib209_cap100_holdoutchems_advi/models/chem_hier_backoff.npz \
  --test-csv output/rt_pymc_chem_hier_bench/holdout_chems/lib209_cap100_seed42_frac0.1/test.csv \
  --backoff-mean-mode chem \
  --max-test-rows 0 --chunk-size 50000 --n-ppc-draws 200 \
  --label holdout_chems --write-by-species-group

poetry run python scripts/pipelines/eval_rt_pymc_posterior_predictive.py \
  --train-output-dir output/tmp_pymc_chem_hier_lib209_cap100_holdoutchems_advi \
  --backoff-npz output/tmp_pymc_chem_hier_lib209_cap100_holdoutchems_advi/models/chem_hier_backoff.npz \
  --test-csv output/rt_pymc_chem_hier_bench/holdout_chems/lib209_cap100_seed42_frac0.1/test.csv \
  --backoff-mean-mode cluster_only \
  --max-test-rows 0 --chunk-size 50000 --n-ppc-draws 200 \
  --label holdout_chems_cluster_only --write-by-species-group

# Plot backoff comparison by species_group_raw
poetry run python scripts/pipelines/compare_rt_models_by_group.py \
  --hier-csv output/tmp_pymc_chem_hier_lib208_cap100_holdoutchems_advi/results/rt_eval_ppc_by_species_group_holdout_chems.csv \
  --lasso-csv output/tmp_pymc_chem_hier_lib208_cap100_holdoutchems_advi/results/rt_eval_ppc_by_species_group_holdout_chems_cluster_only.csv \
  --hier-label "chem_hier backoff" --lasso-label "cluster-only backoff" --label holdout_chems \
  --output docs/models/images/rt_pymc_collapsed_ridge_report/rt_compare_rmse_by_species_group_lib208_holdout_chems.png

poetry run python scripts/pipelines/compare_rt_models_by_group.py \
  --hier-csv output/tmp_pymc_chem_hier_lib209_cap100_holdoutchems_advi/results/rt_eval_ppc_by_species_group_holdout_chems.csv \
  --lasso-csv output/tmp_pymc_chem_hier_lib209_cap100_holdoutchems_advi/results/rt_eval_ppc_by_species_group_holdout_chems_cluster_only.csv \
  --hier-label "chem_hier backoff" --lasso-label "cluster-only backoff" --label holdout_chems \
  --output docs/models/images/rt_pymc_collapsed_ridge_report/rt_compare_rmse_by_species_group_lib209_holdout_chems.png
\end{verbatim}

\section{Next steps}

The partially-collapsed chem\_hier model matches the collapsed ridge accuracy on \emph{seen groups} while producing better-calibrated
intervals and providing a backoff path for \emph{unseen compounds}. The next iteration should focus on (1) making cold-start predictions
useful and (2) bringing training time down:
\begin{itemize}
  \item \textbf{Improve cold-start (unseen groups) beyond intercept-only backoff.}
        The held-out-chem experiment is strict zero-shot and currently reverts unseen-group slopes to $\bm{0}$, which yields poor RMSE.
        Add a shared slope mean (e.g., global and/or cluster-level) so unseen groups can inherit run-covariate response.
  \item \textbf{Re-run held-out evaluations after slope sharing.}
        Repeat the held-out-chem split and report RMSE/coverage (global and by \texttt{species\_group\_raw}) to quantify the lift from
        chemistry and the shared slope head.
  \item \textbf{Tune ADVI for speed/quality.}
        Sweep \texttt{--advi-steps} (and optionally \texttt{--advi-draws}) on cap5 first, then cap100, and record the time vs RMSE/coverage
        trade-off.
  \item \textbf{Optionally learn $\lambda$ once the above is stable.}
        Enable \texttt{--lambda-mode learn} and check whether it improves intervals without harming seen-group RMSE.
  \item \textbf{Keep the benchmark comparison stable.}
        Continue reporting: (a) global metrics on realtest (seen groups only), (b) per-\texttt{species\_group\_raw} RMSE vs lasso
        (including lasso compound coverage), and (c) the held-out-chem experiment for cold-start.
\end{itemize}

\end{document}
